\section{The documents, and what each contains}\label{sec:corpus}

\subsection{The editions and how they relate}

\begin{center}\small
\begin{tabularx}{\linewidth}{@{}>{\raggedright\arraybackslash}p{0.19\linewidth}>{\raggedright\arraybackslash}p{0.23\linewidth}>{\raggedright\arraybackslash}X@{}}
\toprule
Edition & Where & Contents and relation to the others\\
\midrule
Project reader, 488 pages, 7 August 2026 & repository root, with source and Lean archives & The earlier edition. The README says it is preserved and is not the newest reader.\\
Cumulative archive, 619 pages, 31 August 2026 & Zenodo 22666493 and 22678971, file 00 & The 488-page reader continued. Its main source has the same section skeleton, now 60{,}561 lines instead of 57{,}452. Fourteen satellite files add §§10--19 and the carrier refinements §§24.1--24.3. It begins from the source of the earlier Zenodo record 21845035, version~69 (archive p.~3). It is the front file of the Zenodo collection.\\
Bounded congruence transport, 67 pages, 8 September & Zenodo 22666493, files 02--03 & The 58-statement supplement.\\
Expanded supplement, 86 pages, 9 September & Zenodo 22678971, files 07--08; repository \texttt{research/\allowbreak expanded-2026-09-08/} & The same, with 14 further statements: 72 statements in 15 source files, with checkers (§\ref{ssec:supplement}).\\
ES/Fable reader, 124 pages, 20 September & repository \texttt{research/\allowbreak incoming/\allowbreak es-fable-zeta-bridge-\allowbreak 20260920/} & A coordinate bridge from Erdős--Straus solutions to the Fable cover and a weighted conductor of the zeta programme (§\ref{ssec:crosswalk}).\\
Continuation, 17--23 September & repository README, ``Start reading'' (28 items); \texttt{research/\allowbreak incoming/} & 41 dated packages; the results bench R01--R10 and the integration arguments X1--X7 (§\ref{ssec:continuation}).\\
\bottomrule
\end{tabularx}
\end{center}

These editions have different contents; the repository preserves them rather than overwriting them (README, edition table).

\subsection{The 619-page archive, section by section}\label{ssec:archive}

Page numbers are those of the PDF. The column ``Bearing'' says whether a section bears on the equation:
\begin{itemize}
\item \emph{direct}: it proves or tests solvability;
\item \emph{literature}: it audits a published approach to the equation;
\item \emph{side}: a construction not about solvability.
\end{itemize}
The status column refers to this reader.

\begin{center}\small
\begin{longtable}{@{}>{\raggedright\arraybackslash}p{0.06\linewidth}>{\raggedright\arraybackslash}p{0.08\linewidth}>{\raggedright\arraybackslash}p{0.45\linewidth}>{\raggedright\arraybackslash}p{0.13\linewidth}>{\raggedright\arraybackslash}p{0.17\linewidth}@{}}
\toprule
§ & pp. & Content & Bearing & Status here\\
\midrule
\endhead
front & 3--5 & Research status: the census figures of §24, the four families, the record-shell results, and an explicit list of non-claims & summary & read\\
1 & 6--16 & Scope, road map, mathematical status, attribution, lineage & summary & read in part\\
2 & 16--18 & The discriminant-five identity $(2n+1)^2-5=4(n^2+n-1)$ and its root involution, from u/CivQ17 & side & located\\
3 & 18--20 & The shell $p_*=8{,}803{,}369$, $R_*=107$, reached through Busy-Beaver values ($S(4)=107$, $\mathrm{Space}(5)=12289$) & direct (one shell) & $R=107$ is the least occupied shell at $p_*$ (Proposition~\ref{prop:records})\\
4 & 20--22 & Alpöge's constant-Jacobian map $F:\C^3\to\C^3$ with a three-point fibre & side & see \S\ref{sec:overlaps}\\
5 & 22 & Withdrawn constructions (a cube-root interpolation on the Fable fibre), with the point of failure & side & read\\
6 & 22--27 & Central-angle and inversive-circle theorems of u/UmbrellaCorp\_HR; a Lorentz reflection & side & located\\
7 & 27--31 & The complete fibre of $F$ over $(-\tfrac14,0,0)$ and a marked-factor lift & side & see \S\ref{sec:overlaps}\\
8 & 31--38 & The singular $C_1$--$C_2$--$C_3$ support object; dihedral actions & side & located\\
9 & 38--93 & The exact divisor-shell bijection, then directed blade transport and a quarter-radius inversion & direct (§9 opening), side (the rest) & bijection proved here (Theorem~\ref{thm:shell})\\
10 & 93--98 & Residual sieve modulo $840$ and the Jacobi-symbol barrier & direct & proved here (Propositions~\ref{prop:mod840}, \ref{prop:jacobi})\\
11 & 98--100 & Elsholtz--Tao Type~I and Type~II coordinates mapped to the residual shell & literature & located\\
12 & 100--106 & Adjacent programmes (López, Dyachenko, Bradford): what each proves, and where a claimed bridge fails & literature & located\\
13 & 106--114 & Mordell's constructions and Rosati's theorem; Salez's seven charts; Bright--Loughran surfaces; the Vaughan boundary & literature & located\\
14 & 114--119 & López's structural solutions: exact divisor coordinates and a quantifier repair & literature & located\\
15 & 119--124 & The Star--Kneser programme: a one-way forcing theorem and a retracted converse & direct & located\\
16 & 124--127 & The mod-$107$ deficit progression, following u/CommonCareful3149 & direct & progression and its corollary reproduced (§\ref{ssec:records})\\
17 & 127--135 & Two-fraction criteria; Monks--Velingker; Ionascu--Wilson classes; Zelator's prime-square classification & literature & classes replayed (§\ref{ssec:records})\\
18 & 135--148 & The residual-eleven support ideal, its colons, and the paired $R=7/R=11$ affine theorem & direct & located\\
19 & 148--157 & Strict record shells: an open conjecture and the boundary example $p=806{,}521$ & direct & records and deficits reproduced (§\ref{ssec:records})\\
20 & 157--179 & The shell $p=12289$, $R=11$, $a=3075$: counts $(3,6,3)$, its dihedral orbit, and the Star--Kneser quotient certificate & direct (one shell) & counts reproduced: they are $(|E_a|,|M_a|,|E_a|)$, and $a$ is the least occupied shell; the rest located\\
21 & 179--200 & The split-zero semiring $G(R)=R\sqcup\{\tau\}$ and the $C_1$ carrier & side & see \S\ref{sec:overlaps}\\
22 & 200--201 & Rigidified denominator towers over $\Ff_1$ & side & located\\
23 & 201--250 & Mirror completion of the Boolean carrier; the $12289$ orbit matrix & side & located\\
24 & 250--604 & Tomita--Takesaki transport and $C_3$ modular data (opening part); then the carrier census: the $R=167$ and $R=239$ carrier theorems, the refresh of the $R=167$ coefficient, and the return to the $R=23$ lineage (pp.~451--604), including the four families (p.~423) & direct (census), side (operator part) & four families proved here (Theorem~\ref{thm:families}); census located (§\ref{ssec:census})\\
25 & 605--616 & A map from Busy-Beaver extrema to finite shell carriers & side & see \S\ref{sec:overlaps}\\
\bottomrule
\end{longtable}
\end{center}

So about 100 of the 619 pages bear directly on the equation: §§3, 9 (opening), 10, 15, 16 and 18--20, and the census part of §24. The literature audits of §§11--14 and 17 add about 35 pages. The rest is side constructions, whose connection to the equation the archive itself describes as an interface still to be supplied by an explicit typed map (archive p.~4).

\subsection{The bounded-transport supplement (72 statements)}\label{ssec:supplement}

The supplement starts from the primes $p=12h+1$ and all shells $3h+1\le a\le9h$ of Theorem~\ref{thm:shell}. Its sections are:
\begin{enumerate}[label=\arabic*.]
\item original shells, coordinates and integer reconstruction, with the first-two-shell sieve and an all-shell no-hit sieve;
\item local congruences with all integral lifts retained, with an integral cyclic bridge;
\item bounded divisibility and simultaneous congruences: unary Boolean transport, a shared-variable Chinese-remainder construction credited to a colleague, raw-scale hit incidence, and a boundary-swap completion;
\item exact cross-shell transport at a fixed rational ratio;
\item a primitive Euclidean successor with its full arithmetic domain;
\item arithmetic localisation along the original norm subdivision, with a reading of Connes' primitive intersections;
\item a strictly decreasing denominator with exact arithmetic return;
\item the Navier--Stokes auxiliary cover: the torus matrix $J=\begin{pmatrix}3&1\\1&5\end{pmatrix}$ of the announced Navier--Stokes manuscript, its covering fibres and character corrections, and sampling and reconstruction estimates.
\end{enumerate}
The supplement's own README states that its results prove or disprove neither the Erdős--Straus conjecture nor the Riemann hypothesis nor the Yang--Mills mass gap, and that they do not validate the Navier--Stokes construction.

\status{Reproduced in part. A second-pass audit, interrupted by a usage limit before it wrote its report, recomputed the finite content of the statements of Sections 1--7. Its scripts, which are independent of the workbench's code, found every checked statement correct:
\begin{itemize}
\item the gluing and Chinese-remainder theorems on random systems and on all shells of small primes;
\item the ratio-transport domain on all $57{,}882$ ordered shell pairs of the primes $p\equiv1\pmod{12}$ below $260$;
\item the shear and no-two-shears theorems on all shells of $131$ primes below $4000$;
\item the Connes marked successor on all primes $p\equiv1\pmod{12}$ below $3\cdot10^5$;
\item the fixed-$y$ successor on all $3{,}114$ source witnesses below $1500$;
\item the orchard partition on all $42$ admissible triples below $1300$;
\item the boundary four-candidate law on all $72{,}282$ raw records below $400$;
\item the unary box, the input/output incidence and the raw-scale theorems on their stated ranges.
\end{itemize}
One of its checks, on the smallest member of a boundary orbit, reported mismatches. The script expected the pattern $(\min(A,B),\max(A,B))$ in every case. The supplement's own table (\texttt{boundary\_swap\_completion.tex}, lines 296--309) gives a one-point orbit's own point as its smallest member, which is what the script found. The mismatch is therefore in the script, not in the supplement. The same pass also looked at scope, and its search led to Proposition~\ref{prop:twoshears} below. The general proofs were not re-derived, except that of Section~5's no-two-shears theorem; the auxiliary-torus Section 8 was not checked. The scripts are in \note{checks/es\_reader/second\_pass\_partial/}.}

Section 5 of the supplement studies the two unimodular shears $L(m,n)=(m+n,n)$ and $J(m,n)=(m,m+n)$ as maps between adjacent shells. It encodes the solutions with first denominator $a$ by the set
\[
 \mathcal W_a=\{(m,n)\in\ZZ_{>0}^2:\ m\mid pa,\ n\mid pa,\ \gcd(m,n)=1,\ R_a\mid m+n\},
\]
which is in bijection with the ordered solutions $(a,y,z)$ through $y=pak/n$, $z=pak/m$, $k=(m+n)/R_a$ (supplement §1, first lemma; this is Theorem~\ref{thm:shell}(b) in other coordinates). A \emph{shear edge} is a step $(a,m,n)\mapsto(a+1,L(m,n))$ or $(a,m,n)\mapsto(a+1,J(m,n))$ with $(m,n)\in\mathcal W_a$ and the image in $\mathcal W_{a+1}$. The supplement proves, for $p=12h+1$ and shells $3h+1\le a\le9h$, that no two shear edges compose (Theorem ``No two consecutive positive Euclidean shell steps'', \texttt{publication.tex}, line~933). Its proof uses the hypothesis on $p$ at one step only, and there only through $p\equiv1\pmod3$.

\begin{proposition}[the no-two-shears theorem, exact scope]\label{prop:twoshears}
Let $p$ be a prime, and consider all shells $a$ with $p/4<a<p$.
\begin{enumerate}[label=(\alph*)]
\item If $p\not\equiv2\pmod3$, no two shear edges compose: there is no path $(a,m,n)\mapsto(a+1,m',n')\mapsto(a+2,m'',n'')$ of shear edges.
\item For $p\equiv2\pmod3$ such paths exist. At $p=131$ one is $(33,1,33)\mapsto(34,1,34)\mapsto(35,1,35)$, with $R_a=1,5,9$. At $p=1613$, where $p\equiv1\pmod4$, one is $(404,1,404)\mapsto(405,1,405)\mapsto(406,1,406)$, with $R_a=3,7,11$.
\end{enumerate}
\end{proposition}
\status{\textsf{Generalised here}, with the workbench's proof. The supplement states (a) for $p\equiv1\pmod{12}$ and $3h+1\le a\le9h$ only. Checked on the full graph of every prime below $2000$ (script \note{esr\_twoshears.py}: no path for $p\equiv1\pmod3$; $14$ paths for $p\equiv2\pmod3$), and by a complete search up to $p\le10^7$ that uses the shape of an edge proved below: it finds $19{,}558$ paths of two $L$-edges, all at primes $p\equiv2\pmod3$, and as many of two $J$-edges.}
\begin{proof}
\emph{The shape of one edge} (the supplement's classification, whose proof needs only $a+1<p$). Let $(a,m,n)\mapsto(a+1,m+n,n)$ be an $L$-edge. Then $n$ divides $pa$ and $p(a+1)$, whose greatest common divisor is $p$. If $n=p$, then $p\nmid m$ by coprimality, and $m+p$ divides $p(a+1)$ and is prime to $p$, so $m+p\mid a+1$; but $m+p>p>a+1$. So $n=1$. If $p\mid m$, then $m+1$ is prime to $p$ and exceeds $a+1$, and yet divides $p(a+1)$, which is impossible; so $p\nmid m$ and $m\mid a$. Then $m+1\le a+1<p$, so $m+1\mid a+1$. Finally $R_a\mid m+1$ and $R_{a+1}=R_a+4\mid m+2$. For a $J$-edge, exchange $m$ and $n$; the exchange preserves every $\mathcal W_a$.

(a) An $L$-edge ends at $(a+1,m+1,1)$, whose first coordinate is at least $2$; a $J$-edge starts at a pair whose first coordinate is $1$. So an $L$-edge is never followed by a $J$-edge, and by the exchange a $J$-edge is never followed by an $L$-edge. Suppose two $L$-edges pass through $(a+i,m+i,1)$ for $i=0,1,2$. By the shape of each edge, $m+i\mid a+i$ and $R_a+4i\mid m+i+1$ for $i=0,1,2$. Put $K=a-m$. Then $m+i\mid K$ for each $i$, and one of $m,m+1,m+2$ is divisible by $3$, so $3\mid K$. Also $R_a+4i$ divides $4(m+i+1)-(R_a+4i)=4m+4-R_a=:C$, and one of $R_a,R_a+4,R_a+8$ is divisible by $3$, so $3\mid C$. But $C=4m+4-4a+p=p+4-4K\equiv p+1\pmod3$. Hence $p\equiv2\pmod3$. Two $J$-edges are excluded by the exchange.

(b) At $p=131$: $R_{33}=1$ divides $1+33$; $R_{34}=5$ divides $35$; $R_{35}=9$ divides $36$; each coordinate divides $pa$, and each pair is coprime. At $p=1613$ (prime): $3\mid405$, $7\mid406=7\cdot58$ and $11\mid407=11\cdot37$. Each step is $J(1,n)=(1,n+1)$.
\end{proof}
So the hypothesis $p\equiv1\pmod{12}$ of the supplement can be weakened to $p\not\equiv2\pmod3$, and not further. The primes $p\equiv2\pmod3$ are solvable in any case (Section~\ref{sec:core}); the proposition concerns the transport graph, not solvability.

\subsection{The 124-page ES/Fable reader}\label{ssec:crosswalk}

It encodes a solution of $4/p=1/x+1/y+1/z$ as the binary quartic
\[
H_{\mathrm{ES}}(U,V)=-S^{-1}(U-pV)(U-xV)(U-yV)(U-zV),\qquad S=p+x+y+z,
\]
with normalised coefficients $u_0=-1/S$, $u_2=-\bigl(p(x+y+z)+xy+yz+zx\bigr)/S$, $u_3=5xyz/S$ and $u_4=-pxyz/S$. It then transports the quartic to the Fable cover and to a weighted conductor of the zeta programme.

\begin{proposition}\label{prop:quartic}
On solutions, $p=-5u_4/u_3$ and
\[
625u_0u_4^3-125u_3u_4^2+25u_2u_3^2u_4-4u_3^4=0 .
\]
\end{proposition}
\status{Proved here; the README of the package states both, and the relation was also checked in \note{21\_} (Proposition~21.6).}
\begin{proof}
Write $\sigma_1=x+y+z$, $\sigma_2=xy+yz+zx$ and $P=xyz$. The equation is $4P=p\sigma_2$, so the third elementary symmetric function of $p,x,y,z$ is $p\sigma_2+P=5P$; this gives $u_3$, and $p=-5u_4/u_3$ follows. Substituting and multiplying by $S^4/625$, the left side becomes $P^3\bigl(p^3-p^2S+p(p\sigma_1+\sigma_2)-4P\bigr)=P^3(p\sigma_2-4P)=0$.
\end{proof}
The relation is thus the equation itself, written in these coordinates. The Fable and conductor maps built on it are \textsf{located}; their comparison with the zeta programme is in the zeta audit (\note{47\_}, overlap OV-13).

\subsection{The continuation of 17--23 September}\label{ssec:continuation}

The README's ``Start reading'' list has 28 items with proof packages in \texttt{research/incoming/}. The results bench R01--R10 and the integration arguments X1--X7 accompany them.
\begin{itemize}
\item The first-pass audit \note{43a\_} lists them all.
\item Section~\ref{sec:items} states the eleven items audited there.
\item Section~\ref{sec:negative} gives their negative results.
\end{itemize}
